\documentclass[12pt,a4paper]{article}
\usepackage[utf8]{inputenc}
\usepackage[T1]{fontenc}
\usepackage[brazilian]{babel}
\usepackage[margin=2.5cm]{geometry}
\usepackage{listings}
\usepackage{xcolor}
\usepackage{booktabs}
\usepackage{hyperref}

\lstset{
  basicstyle=\ttfamily\footnotesize,
  breaklines=true,
  frame=single,
  backgroundcolor=\color{gray!8},
  literate={ã}{{\~a}}1 {á}{{\'a}}1 {é}{{\'e}}1 {ó}{{\'o}}1 {ú}{{\'u}}1 {ç}{{\c c}}1 {í}{{\'i}}1,
  keepspaces=true
}

\begin{document}

\begin{titlepage}
  \centering
  \vspace*{3cm}
  {\Large Teste de Software\par}
  \vspace{1.5cm}
  {\LARGE\bfseries Detecção de Bad Smells e Refatoração Segura\par}
  \vspace{3cm}
  {\large Miguel Figueiredo Diniz\par}
  {\large Matrícula: 853959\par}
  \vfill
  {\small \url{https://github.com/DevMiguelDiniz/bad-smells-js-refactoring}\par}
  \vspace{1cm}
\end{titlepage}

\section{Análise de Smells (análise manual)}
Lendo \texttt{src/ReportGenerator.js} identifiquei três Bad Smells principais.

\paragraph{1) Long Method / complexidade ciclomática alta.} O método \texttt{generateReport} concentra cabeçalho, filtragem por perfil, cálculo de prioridade, formatação das linhas, soma do total e rodapé, com \texttt{for} e \texttt{if/else} aninhados em até 4 níveis. Na manutenção, qualquer alteração exige entender o método inteiro. Nos testes, o número de caminhos (tipo de relatório $\times$ perfil $\times$ valor do item) cresce de forma multiplicativa, e as unidades de lógica não podem ser testadas isoladamente.

\paragraph{2) Código duplicado.} As linhas CSV e HTML e o \texttt{total += item.value} aparecem repetidos nos ramos ADMIN e USER, e o teste do tipo de relatório (CSV/HTML) se repete no cabeçalho, no corpo (duas vezes) e no rodapé. Uma mudança de formato precisa ser feita em vários lugares, com risco de divergência, e cada cópia exige cobertura própria.

\paragraph{3) Condicionais complexas e números mágicos.} As cadeias de \texttt{if/else} por tipo de relatório e por perfil violam o princípio aberto/fechado: um novo formato (ex.: JSON) ou perfil obriga a editar o método e re-testar tudo. Os literais \texttt{1000} (limite de prioridade), \texttt{500} (limite do usuário comum) e as strings \texttt{'ADMIN'}, \texttt{'USER'}, \texttt{'CSV'}, \texttt{'HTML'} não têm nome, então a intenção é implícita. Há ainda um efeito colateral: o método \emph{muta} o item recebido (\texttt{item.priority = true}).

\section{Relatório da Ferramenta (ESLint + SonarJS)}
Configuração: \texttt{.eslintrc.json} com \texttt{eslint:recommended} e \texttt{plugin:sonarjs/recommended}. O repositório usa o limite de complexidade cognitiva \textbf{5} (o enunciado sugere 15). Saída de \texttt{npx eslint src/} \textbf{antes} da refatoração (\texttt{relatorio/eslint\_antes.txt}):

\begin{lstlisting}
src\ReportGenerator.js
  11:3   error  Refactor this function to reduce its Cognitive
         Complexity from 27 to the 5 allowed  sonarjs/cognitive-complexity
  43:14  error  Merge this if statement with the nested one
         sonarjs/no-collapsible-if

2 problems (2 errors, 0 warnings)
\end{lstlisting}

O SonarJS quantificou o que na leitura manual era apenas uma impressão: a complexidade cognitiva é \textbf{27} contra um limite de 5, ponderando o aninhamento (cada nível extra pesa mais), algo difícil de estimar a olho. Também apontou um \texttt{if} colapsável (linha 43) que passa despercebido por parecer uma regra de negócio legítima. Por outro lado, a ferramenta \textbf{não} detectou números mágicos nem a mutação do parâmetro, achados apenas na análise manual. As duas abordagens são complementares.

\section{Processo de Refatoração}
O smell mais crítico é o Long Method com condicionais aninhadas e duplicação (complexidade 27).

\paragraph{Antes} (trecho do corpo):
\begin{lstlisting}[language=JavaScript]
for (const item of items) {
  if (user.role === 'ADMIN') {
    if (item.value > 1000) { item.priority = true; }
    if (reportType === 'CSV') {
      report += `${item.id},${item.name},...`;
      total += item.value;
    } else if (reportType === 'HTML') {
      const style = item.priority ? 'style="font-weight:bold;"' : '';
      report += `<tr ${style}><td>...</td></tr>\n`;
      total += item.value;
    }
  } else if (user.role === 'USER') {
    if (item.value <= 500) {
      if (reportType === 'CSV') { /* ...duplicado... */ }
      else if (reportType === 'HTML') { /* ...duplicado... */ }
    }
  }
}
\end{lstlisting}

\paragraph{Depois} (\texttt{src/ReportGenerator.refactored.js}):
\begin{lstlisting}[language=JavaScript]
generateReport(reportType, user, items) {
  const formatter = formatters[reportType] ?? emptyFormatter;
  const visibleItems = this.selectVisibleItems(user, items);

  this.markPriorityItems(user, visibleItems);

  const total = visibleItems.reduce((sum, item) => sum + item.value, 0);
  const rows = visibleItems.map((item) => formatter.row(item, user));

  return [formatter.header(user), ...rows, formatter.footer(total)]
    .join('').trim();
}
\end{lstlisting}

Técnicas aplicadas:
\begin{itemize}
  \item \textbf{Extract Method:} \texttt{selectVisibleItems} e \texttt{markPriorityItems} isolam filtragem e regra de prioridade.
  \item \textbf{Replace Conditional with Polymorphism (Strategy):} \texttt{csvFormatter} e \texttt{htmlFormatter} expõem a mesma interface (\texttt{header}, \texttt{row}, \texttt{footer}); o \texttt{if/else} por tipo virou uma busca em tabela. Novo formato = novo objeto, sem editar o método.
  \item \textbf{Decompose Conditional / tabela de regras:} \texttt{visibilityRules} mapeia perfil $\to$ predicado.
  \item \textbf{Replace Magic Number with Symbolic Constant:} \texttt{PRIORITY\_THRESHOLD}, \texttt{USER\_MAX\_VALUE}, \texttt{ROLE\_ADMIN}, \texttt{ROLE\_USER}.
  \item \textbf{Eliminação de duplicação:} formatação e soma do total ocorrem uma única vez.
\end{itemize}
O comportamento para entradas inesperadas foi preservado (tipo desconhecido gera relatório vazio; perfil desconhecido gera relatório sem itens).

\section{Validação}
O arquivo \texttt{tests/ReportGenerator.refactored.test.js} é cópia do teste original, alterando \textbf{apenas} o import. Resultado de \texttt{npm test}:

\begin{center}
\begin{tabular}{@{}lc@{}}
\toprule
Suíte & Resultado \\
\midrule
\texttt{ReportGenerator.refactored.test.js} & \textbf{5/5 passam} \\
\texttt{ReportGenerator.test.js} (original) & 4/5 passam \\
\bottomrule
\end{tabular}
\end{center}

\texttt{npx eslint src/ReportGenerator.refactored.js}: \textbf{0 erros} (antes: 2 erros, complexidade 27).

\paragraph{Observação importante.} O teste original \emph{já falhava antes de qualquer alteração} (``HTML completo para Admin''). O código original gera \texttt{<tr >} (com espaço, vindo do template quando \texttt{style} é vazio), mas o teste espera \texttt{<tr><td>1</td>...}. É um bug latente do código original, não uma regressão. A versão refatorada gera \texttt{<tr>} e \texttt{<tr style="font-weight:bold;">}, passando em todos os testes. Conforme a regra do trabalho, o arquivo e o teste originais não foram modificados.

\section{Conclusão}
A suíte de testes funcionou como rede de segurança: validou o comportamento externo (saída CSV/HTML, filtros, totais) e permitiu reescrever a estrutura interna com confiança, além de revelar um bug preexistente (\texttt{<tr >}) que a leitura do código não evidenciava. A redução dos smells (complexidade cognitiva 27 $\to$ abaixo do limite, duplicação removida, constantes nomeadas) torna o código mais legível, mais fácil de estender (novos formatos e perfis) e mais testável, pois cada parte (formatter, regra de visibilidade, prioridade) pode ser testada isoladamente.

\end{document}
